\chapter{Weyl 变换与 Wigner 函数}
\label{ch:weyl-wigner}

\begin{learningobjectives}
完成本章后，读者应能：
\begin{enumerate}
  \item 从密度矩阵写出 Wigner 函数，并验证归一化和边缘分布；
  \item 计算简单算符的 Weyl 符号；
  \item 解释为什么 Wigner 平均对应对称排序；
  \item 正确抽样相干态和真空态的单模 Wigner 分布。
\end{enumerate}
\end{learningobjectives}

\section{从密度矩阵到实函数}

先从有量纲的位置--动量对 $\hat x,\hat p_x$ 出发，满足
\begin{equation}
  \comm{\hat x}{\hat p_x}=\ii\hbar.
\end{equation}
密度算符 $\rhohat$ 在坐标表象中的矩阵元 $\mel{x_1}{\rhohat}{x_2}$ 依赖两个坐标。引入中心坐标 $x=(x_1+x_2)/2$ 和相对坐标 $y=x_1-x_2$，再对 $y$ 作 Fourier 变换，得到 Wigner 函数
\begin{equation}
  W(x,p_x)=\frac{1}{2\pi\hbar}
  \int_{-\infty}^{\infty}\dd y\,
  \ee^{-\ii p_x y/\hbar}
  \mel{x+y/2}{\rhohat}{x-y/2}.
  \label{eq:wigner-coordinate-definition}
\end{equation}
这是 Wigner 准分布的标准坐标表象定义；不同归一化与量子光学中的常用复
振幅约定可与原始工作及综述交叉核对\cite{Wigner1932,Hillery1984}。
这里选择的系数使
\begin{equation}
  \int\dd x\,\dd p_x\,W(x,p_x)=\Tr\rhohat=1.
  \label{eq:wigner-normalization}
\end{equation}

\begin{derivationbox}
对\cref{eq:wigner-coordinate-definition} 先积分 $p_x$，利用
\[
  \int\dd p_x\,\ee^{-\ii p_x y/\hbar}=2\pi\hbar\,\delta(y),
\]
便得到 $\int\dd p_x\,W(x,p_x)=\mel{x}{\rhohat}{x}$。再积分 $x$ 即为迹。这个计算也说明 Wigner 函数的位置边缘分布是真正的概率密度。
\end{derivationbox}

Hermitian 性 $\rhohat^\dagger=\rhohat$ 保证 $W$ 为实函数。实函数并不意味着正函数；允许负值正是它能够同时编码不对易量信息的代价。另一方面，负值不是“量子性的必要条件”，真空态、相干态、热态和所有单模高斯态都可以具有处处非负的 Wigner 函数。

\section{边缘分布}

直接积分\cref{eq:wigner-coordinate-definition} 得
\begin{align}
  \int\dd p_x\,W(x,p_x)&=\mel{x}{\rhohat}{x},\\
  \int\dd x\,W(x,p_x)&=\mel{p_x}{\rhohat}{p_x}.
  \label{eq:wigner-marginals}
\end{align}
不过，\cref{eq:wigner-marginals} 不能解释为一次实验同时测得了确定的 $x$ 和 $p_x$；它只说明同一个准分布能够产生不同测量设置的边缘统计。有量纲的 $x$ 与 $p_x$ 不能在未选定尺度时直接相加；旋转正交分量将在下面引入无量纲变量后定义。

\section{Weyl 符号}

任意算符 $\hat A$ 的 Weyl 符号定义为
\begin{equation}
  A_W(x,p_x)=\int_{-\infty}^{\infty}\dd y\,
  \ee^{-\ii p_x y/\hbar}
  \mel{x+y/2}{\hat A}{x-y/2}.
  \label{eq:weyl-symbol-definition}
\end{equation}
密度算符的 Weyl 符号与 Wigner 函数相差归一化因子。两者配对给出迹公式
\begin{equation}
  \Tr(\rhohat\hat A)
  =\int\dd x\,\dd p_x\,W(x,p_x)A_W(x,p_x).
  \label{eq:weyl-trace-rule}
\end{equation}

Weyl 映射把完全对称排序的算符积映射为普通乘积。例如
\begin{align}
  (\hat x)_W&=x, & (\hat p_x)_W&=p_x,\\
  \left(\frac{\hat x\hat p_x+\hat p_x\hat x}{2}\right)_W&=xp_x,
  & (\hat x^2)_W&=x^2.
\end{align}
但一般情况下 $(\hat A\hat B)_W\neq A_WB_W$。算符乘积对应的星积将在第3章推导。

\section{玻色模式的复相空间}

对单个玻色模式，另行定义无量纲正交分量
\begin{equation}
  \ha=\frac{\hat Q+\ii\hat P}{\sqrt{2}},
  \qquad
  \alpha=\frac{Q+\ii P}{\sqrt{2}},
  \qquad \comm{\hat Q}{\hat P}=\ii,
  \qquad \comm{\ha}{\had}=1.
  \label{eq:mode-quadratures}
\end{equation}
物理系统若给出特征长度 $x_0$，一种常见选择是 $\hat Q=\hat x/x_0$、$\hat P=x_0\hat p_x/\hbar$；不同选择只是辛尺度变换，不能把两套变量的测度混用。复平面面积元为 $\dd^2\alpha=\dd(\Re\alpha)\dd(\Im\alpha)=\dd Q\dd P/2$。本书在复振幅表象中把 Wigner 函数归一化为
\begin{equation}
  \int\dd^2\alpha\,W(\alpha,\alpha^*)=1.
\end{equation}

现在可以定义可测的旋转正交分量
\begin{equation}
  \hat X_\theta=\hat Q\cos\theta+\hat P\sin\theta.
\end{equation}
对 Wigner 函数沿垂直于 $X_\theta$ 的方向作线积分，就得到该正交分量的概率分布；这也是量子层析中 Radon 反演的基础。

另一种常用定义从对称特征函数开始：
\begin{align}
  \chi_W(\lambda)&=\Tr\!\left[\rhohat\,
  \ee^{\lambda\had-\lambda^*\ha}\right],\\
  W(\alpha,\alpha^*)&=\frac{1}{\pi^2}
  \int\dd^2\lambda\,\chi_W(\lambda)
  \ee^{\lambda^*\alpha-\lambda\alpha^*}.
  \label{eq:wigner-characteristic}
\end{align}
指数中的位移算符同时包含 $\ha$ 和 $\had$，这正是对称排序出现的来源。

\section{真空态与相干态}

相干态满足 $\ha\ket{\alpha_0}=\alpha_0\ket{\alpha_0}$，其 Wigner 函数为
\begin{equation}
  W_{\alpha_0}(\alpha)
  =\frac{2}{\pi}\exp\!\left[-2|\alpha-\alpha_0|^2\right].
  \label{eq:coherent-wigner}
\end{equation}
真空态是 $\alpha_0=0$ 的特例。对\cref{eq:coherent-wigner} 的直接抽样为
\begin{equation}
  \alpha=\alpha_0+\frac{\xi_1+\ii\xi_2}{2},
  \qquad \xi_1,\xi_2\sim\mathcal{N}(0,1),
  \label{eq:coherent-sampling}
\end{equation}
其中两个实高斯变量相互独立。因此
\begin{equation}
  \ensavg{\alpha}=\alpha_0,
  \qquad
  \ensavg{|\alpha-\alpha_0|^2}=\frac12.
  \label{eq:vacuum-half-quantum}
\end{equation}
这个“半个量子”是正交分量真空噪声的相空间表示。

\begin{numericalbox}
抽样程序的第一项单元测试应检查\cref{eq:vacuum-half-quantum}，而不是只检查随机数看起来是否像高斯。若复噪声误写成 $(\xi_1+\ii\xi_2)/\sqrt{2}$，每个模式会被加入一个而不是半个量子的对称排序噪声。
\end{numericalbox}

\section{排序修正：粒子数的例子}

由 $\comm{\ha}{\had}=1$，
\begin{equation}
  \frac{\had\ha+\ha\had}{2}=\had\ha+\frac12.
\end{equation}
左侧是对称排序形式，其 Weyl 符号为 $|\alpha|^2$，所以
\begin{equation}
  (\had\ha)_W=|\alpha|^2-\frac12,
  \qquad
  \qavg{\had\ha}=\ensavg{|\alpha|^2}-\frac12.
  \label{eq:number-weyl-symbol}
\end{equation}
对相干态使用\cref{eq:coherent-sampling}，可得 $\ensavg{|\alpha|^2}=|\alpha_0|^2+1/2$，从而恢复 $\qavg{\had\ha}=|\alpha_0|^2$。

更高阶正规排序矩需要更多收缩项。例如
\begin{equation}
  (\had\had\ha\ha)_W
  =|\alpha|^4-2|\alpha|^2+\frac12.
  \label{eq:quartic-weyl-symbol}
\end{equation}
这类修正在接触相互作用和二阶关联函数中不可忽略。

\section{Wigner 函数能告诉我们什么}

归一化、边缘分布和矩使 Wigner 函数成为连接状态与实验统计的工具。它还满足纯度关系
\begin{equation}
  \Tr(\rhohat^2)=2\pi\hbar\int\dd x\,\dd p_x\,W(x,p_x)^2.
  \label{eq:wigner-purity}
\end{equation}
但不能仅凭一幅 Wigner 图断言动力学机制。负值可以见证某些非高斯非经典性；处处为正则不代表状态可以由经典统计完全解释。对多模场，完整 Wigner 函数存在于高维空间，实际计算主要依赖抽样和低阶边缘量，而不是绘制整个函数。

\begin{warningbox}
文献中的 Wigner 函数可能使用不同的面积元、$\pi$ 因子和正交分量归一化。引用抽样公式前必须同时核对：使用的是有量纲 $(x,p_x)$ 还是无量纲 $(Q,P)$、对易关系、$\alpha$ 的定义、积分测度和真空方差。
\end{warningbox}

\section*{本章小结}

Wigner 函数是密度矩阵相对坐标的 Fourier 变换，能够正确给出正交分量边缘分布。Weyl 符号把算符期望转化为相空间积分，并天然对应完全对称排序。相干态 Wigner 函数的真空宽度产生每模半量子噪声，正规排序观测量必须再作相应修正。

\begin{exercises}
  \item 从\cref{eq:wigner-coordinate-definition} 证明动量边缘分布，即\cref{eq:wigner-marginals} 的第二式。
  \item 验证\cref{eq:coherent-wigner} 在 $\dd^2\alpha$ 测度下归一化为 $1$。
  \item 使用对称排序推导\cref{eq:number-weyl-symbol}。
  \item 对真空态计算 $\ensavg{|\alpha|^4}$，再用\cref{eq:quartic-weyl-symbol} 验证 $\qavg{\had\had\ha\ha}=0$。
  \item 证明纯态的 Wigner 函数满足\cref{eq:wigner-purity} 右侧等于 $1$。
\end{exercises}
